\section{Discussion}\label{discussion}

\subsection{When Switching Helps versus When It
Hurts}\label{when-switching-helps-versus-when-it-hurts}

The results paint a clear picture: topology switching pays off when
tasks genuinely change structural character, and is neutral or slightly
harmful when they don't. The 85.7\% switching precision is encouraging,
but the two negative switches reveal a real failure mode. Both occurred
because early execution patterns (a clarification question, a test
failure on the first attempt) looked like iterativeness to the D-I-T
estimator, when the underlying task was actually decomposable. The
estimator needs more than 8 actions of context to distinguish genuine
iterative structure from transient hiccups.

One possible fix is to require higher confidence before early switches.
A phase boundary detected in the first 10 actions could face a stricter
threshold (\(\delta = 0.6\) instead of \(\delta = 0.4\)). We did not
pursue this in the current work to keep the system simple, but it is an
obvious next step.

\subsection{Connection to Prior Adaptive
Work}\label{connection-to-prior-adaptive-work}

MDAgents \citep{kim_2024} demonstrated that adapting orchestration
structure (solo vs.~group) based on task complexity improves medical
decision-making. Our work extends this principle in two ways. First, we
adapt between three structurally distinct topologies rather than two
modes. Second, we adapt during execution rather than at dispatch. The
MDAgents complexity assessment happens once before the agents start
working; our D-I-T estimation continues throughout.

The concurrent AdaptOrch framework from Yu \citep{adaptorch_2026}
provides the closest comparison. Yu showed that topology choice
dominates system performance when LLM capabilities converge across
providers, a finding that aligns with our OrchestraBench results. But
Yu's system selects topology per-task, not per-phase. On the 42
multi-phase tasks in our suite, this distinction matters: static-oracle
(which mirrors Yu's dispatch-time selection) achieves 88.1\% while our
runtime switching reaches 83.3\%. The 4.8-point gap comes from switching
overhead and imperfect phase detection. The 23.8-point gap between our
system and static-debate shows the value of the approach even with its
current limitations.

RouteLLM \citep{routellm_2025} and bandit-based LLM routing
\citep{banditllm_2025} established that routing decisions can be made
with contextual bandits. We apply the same paradigm at a different
granularity: routing execution phases to topologies rather than queries
to models. The bandit formulation transfers cleanly; the challenge is
defining the context (D-I-T features from traces rather than query
embeddings) and managing the longer feedback loops.

\subsection{Limitations}\label{limitations}

We are direct about what this work does not show.

\textbf{Single backbone.} All experiments use Claude Opus 4.6.
Topology-task alignment patterns might differ with other models. Our
prior OrchestraBench work found that secondary backbone results
converged on easy tasks but we did not test switching behavior across
backbones.

\textbf{82 tasks.} The suite is small by machine learning standards.
With 70 hard tasks and 42 multi-phase tasks, our statistical power is
limited. The 78.6\% success rate has a 95\% confidence interval of
roughly $\pm$9.5 points (binomial). Larger-scale evaluation is needed
to confirm the effect size.

\textbf{Phase boundary detection is brittle.} Our Euclidean distance
heuristic with a fixed threshold is simple but not optimal. It misses
gradual transitions and can be fooled by sudden but temporary shifts.
Learned change-point detectors (e.g., BOCPD) would likely perform
better.

\textbf{Switching overhead is non-trivial.} 800 tokens per switch and
8.2\% total overhead means the approach has negative expected value on
tasks where switching doesn't improve outcomes. The cost model helps,
but it depends on accurate reward predictions from the bandit, which are
themselves uncertain early in execution.

\textbf{No cross-task learning.} Each task starts with the same
warm-started bandit. The system does not learn from previous tasks to
improve switching decisions on new ones. A meta-learning extension could
amortize the exploration cost.

\subsection{Future Directions}\label{future-directions}

Three directions seem promising. First, replacing the fixed threshold
with a learned phase boundary detector trained on our annotated phase
boundaries. Second, extending to a continuous topology space where the
system can interpolate between topologies (e.g., a partially
hierarchical debate) rather than switching discretely between three
options. Third, cross-task transfer of switching policies, where the
bandit's posterior from completed tasks initializes the prior for new
tasks of similar type.
